\begin{corollary}[Unique Predecessor Characterization Away from $1$]
\label{cor:unique-predecessor-characterization-away-from-one}
Let $(P,S,1)$ be a Peano system. For every $x \in P$, $x \neq 1$ if and
only if there exists a unique $u \in P$ such that $S(u)=x$.
\hyperref[prf:unique-predecessor-characterization-away-from-one]{Go to proof.}
\end{corollary}
\end{corollarybox}
\LeanFormalizes{cor:unique-predecessor-characterization-away-from-one}{lra-lean}{LRA.VolumeII.PeanoSystems.BasicTheorems}{UniquePredecessorCharacterizationAwayFromOne}{pending}

\begin{remark*}[Standard quantified statement]
For a Peano system $(P,S,1)$,
\[
\forall x \in P\;\bigl(x\neq 1 \Longleftrightarrow
\exists! u \in P\;(S(u)=x)\bigr).
\]
\end{remark*}

\begin{remark*}[Predicate reading]
\[
\forall x \in P\;\bigl(x\neq 1 \Longleftrightarrow
\exists! u \in P\;(S(u)=x)\bigr).
\]
\end{remark*}

\begin{remark*}[Negated quantified statement]
\[
\begin{aligned}
&\text{Let } (P,S,1) \text{ be a Peano system.}\\
&\exists x \in P\;\Bigl(
\bigl(x\neq 1 \land
(\forall u \in P\;(S(u)\neq x)\\
&\qquad\qquad\qquad\lor
\exists u,v \in P\;(S(u)=x \land S(v)=x \land u\neq v))\bigr)\\
&\qquad\qquad\lor
\bigl(x=1 \land \exists u \in P\;(S(u)=x \land
\forall v \in P\;(S(v)=x \Longrightarrow v=u))\bigr)
\Bigr).
\end{aligned}
\]
\end{remark*}

\begin{remark*}[Failure modes]
\begin{description}
  \item[Exposition.]
  The biconditional can fail in two ways: a non-initial element can fail to
  have a unique predecessor, or the distinguished first element can have a
  unique predecessor.
  \item[Characterization failure.]
  \[
  \neg\forall x \in P\;\bigl(x\neq 1 \Longleftrightarrow
  \exists! u \in P\;(S(u)=x)\bigr).
  \]
  \[
  \begin{aligned}
  &\exists x \in P\;\Bigl(
  \underbrace{\bigl(x\neq 1 \land
  (\forall u \in P\;(S(u)\neq x)
  \lor
  \exists u,v \in P\;(S(u)=x \land S(v)=x \land u\neq v))\bigr)}_
  {\text{non-initial element lacks a unique predecessor}}\\
  &\qquad\qquad\lor
  \underbrace{\bigl(x=1 \land \exists u \in P\;(S(u)=x \land
  \forall v \in P\;(S(v)=x \Longrightarrow v=u))\bigr)}_
  {\text{the first element has a unique predecessor}}
  \Bigr).
  \end{aligned}
  \]
\end{description}
\end{remark*}

\begin{remark*}[Interpretation]
Being different from the distinguished first element is exactly the same as
having a unique predecessor. This turns the informal phrase "away from $1$"
into an internal characterization of the successor structure.
\end{remark*}

\begin{remark*}[Source comparison]
This corollary packages the predecessor analysis surrounding Feferman's
Theorem 3.3. The forward direction is the predecessor
existence-and-uniqueness result, while the reverse direction uses the axiom
that $1$ is not a successor. \citet{FefermanNumberSystems1964}
\end{remark*}

\begin{dependencies}
\begin{itemize}
  \item \hyperref[def:peano-system]{Peano System}
  \item \hyperref[ax:peano-base-not-successor]{Distinguished Element Is Not a Successor}
  \item \hyperref[cor:predecessor-exists-unique-away-from-one]{Predecessor Existence and Uniqueness Away from $1$}
\end{itemize}
\end{dependencies}

